\documentclass[10pt,a4paper]{amsart}
\usepackage[margin=19mm]{geometry}
\usepackage{amsmath,amssymb,mathtools}
\usepackage[scr=rsfs,cal=euler]{mathalfa}
\usepackage{xfrac}
\usepackage[unicode,hidelinks,bookmarks=false]{hyperref}
\newcommand{\ie}{i.e.\ }
\newcommand{\cf}{cf.\ }
\newcommand{\resp}{respectively\ }
\newcommand{\Subsection}{\subsection}
\providecommand{\nobreakdash}{\nobreak\hbox{-}}
\setlength{\emergencystretch}{2em}
\multlinegap0pt
\usepackage{bm}
\usepackage{bbm}
\usepackage{dsfont}
\usepackage{xspace}
\usepackage{cancel}
\usepackage{mathtools}
\usepackage{amsthm}
\makeatletter
\@ifundefined{theorem}{\newtheorem{theorem}{Theorem}}{}
\@ifundefined{proposition}{\newtheorem{proposition}[theorem]{Proposition}}{}
\makeatother
\title[Density Spectra of Topological Groups]{Density Spectra of Topological Groups}
\author{Dekui Peng, Yi Zhou}
\date{Source: arXiv:2507.13031v5 (CC BY 4.0)}
\begin{document}
\maketitle

\noindent\emph{Source and licence.} Density Spectra of Topological Groups. Version arXiv:2507.13031v5 (CC BY 4.0), distributed under the Creative Commons Attribution 4.0 International licence, as recorded in the arXiv licence metadata for this exact version and shown on the abstract page https://arxiv.org/abs/2507.13031v5. Full rights notice: licenses/arxiv-2507.13031-CC-BY-4.0.txt.\par\smallskip

\noindent\textit{Editorial scope.} Counted unit: the Proposition whose source label is \texttt{Prop:t} in the article (source0.tex lines 1008--1010, source label \texttt{Prop:t}), delivered together with the article's own complete original proof of it (source0.tex lines 1012--1046, one \texttt{proof} environment of 35 lines). No mathematics is written, repaired or re-derived: statement and proof are the article's own text line for line. Required context retained uncounted: none. Every definition, display and label of those lines is retained unchanged. Omitted from this package and not redistributed: 1--1007, 1011--1011, 1047--1390. Nothing inside a retained excerpt is omitted or altered: every retained body is delivered exactly as the article's lines are written, so no omission receipt of any kind accompanies this package. Citations: the retained excerpts carry no citation command at all, so no bibliography entry of the article is transcribed and no reference is redistributed. Reference adaptation: none was needed and none was made; every reference inside the delivered unit resolves inside it, so no unresolved reference or citation is produced. Printed slips kept literal: no printed word, symbol, hypothesis or number of the counted statement or of its proof was altered. Layout and class: the article's own class is \texttt{amsart} with packages including amsfonts, amsmath, amssymb, eucal, hyperref; the wrapper is the lane's pinned \texttt{amsart} editorial wrapper at 10pt on A4, which restates the theorem-like environments the retained text needs and adds the article's own macro definitions that the retained text needs (none), with extra wrapper packages (bm, bbm, dsfont, xspace, cancel, mathtools). The delivered numbering is the unit's own. Assets: no retained span contains a figure environment, a graphics-inclusion command, a PDF-page inclusion, a table or a TikZ picture; no image file is copied into this package, the package declares no asset dependency and no image file is redistributed with it. Licence: version arXiv:2507.13031v5 of this article is distributed under the Creative Commons Attribution 4.0 International licence, as recorded in the arXiv licence metadata for this exact version and shown on the abstract page https://arxiv.org/abs/2507.13031v5. A full rights notice is delivered at licenses/arxiv-2507.13031-CC-BY-4.0.txt. Characters: every retained excerpt is pure ASCII, so the delivered engine, XeLaTeX with its default Latin Modern text font, reports no missing character.
\begin{proposition}\label{Prop:t}
Let $G$ be a topological group with density $\kappa > \omega$. Then for every regular cardinal $\tau$ such that $t(G) < \tau \le \kappa$, $G$ contains a closed subgroup of density $\tau$.
\end{proposition}

\begin{proof}
We recursively construct a strictly increasing sequence
$\{H_\alpha:\alpha\leq\tau\}$ of closed subgroups of $G$. Choose any
$x_0\in G$ and put $H_0=\overline{\langle x_0\rangle}$. Suppose that
$H_\beta$ and $x_{\beta+1}$ have been chosen for all $\beta<\alpha$, and set
\[
X_\alpha=\{x_0\}\cup\{x_{\beta+1}:\beta<\alpha\}.
\]
Inductively we shall have $H_\alpha=\overline{\langle X_\alpha\rangle}$.
If $\alpha<\tau$, then
\[
d(H_\alpha)\leq |\langle X_\alpha\rangle|
\leq\max\{|\alpha|,\omega\}<\tau\leq\kappa=d(G),
\]
so $H_\alpha\neq G$. Choose $x_{\alpha+1}\in G\setminus H_\alpha$ and put
\[
H_{\alpha+1}=\overline{\langle H_\alpha\cup\{x_{\alpha+1}\}\rangle}.
\]
At a non-zero limit ordinal $\alpha\leq\tau$, put
$H_\alpha=\overline{\bigcup_{\beta<\alpha}H_\beta}$. Since
$\langle X_\alpha\rangle=\bigcup_{\beta<\alpha}\langle X_\beta\rangle$
at limit stages, the identity
$H_\alpha=\overline{\langle X_\alpha\rangle}$ follows by induction.
In particular, $\langle X_\tau\rangle$ is dense in $H_\tau$ and has
cardinality at most $\tau$, so $d(H_\tau)\leq\tau$.

Suppose that $d(H_\tau)<\tau$. Choose a dense subset $A\subseteq H_\tau$ with $|A|<\tau$. Since $\bigcup_{\alpha<\tau}H_\alpha$ is dense in $H_\tau$, for each $a\in A$ there is a set
$B_a\subseteq\bigcup_{\alpha<\tau}H_\alpha$ such that
$a\in\overline{B_a}$ and $|B_a|\leq t(G)<\tau$. Put $B=\bigcup_{a\in A}B_a$. By regularity of $\tau$, $|B|<\tau$. Again using regularity, there is some $\alpha<\tau$ such that $B\subseteq H_\alpha$. Therefore
\[
A\subseteq\overline B\subseteq H_\alpha,
\]
and hence $H_\tau=\overline A\subseteq H_\alpha$, contradicting
$H_\alpha\subsetneq H_{\alpha+1}\subseteq H_\tau$. Thus $d(H_\tau)=\tau$.
\end{proof}

\end{document}
